\documentclass[10pt]{article}
\usepackage[margin=20mm]{geometry}
\usepackage{amsmath,amssymb,fontspec,unicode-math,xurl,hyperref}
\IfFontExistsTF{DejaVu Serif}{\setmainfont{DejaVu Serif}}{\setmainfont{Latin Modern Roman}}
\setmathfont{latinmodern-math.otf}
\setlength{\parindent}{0pt}\setlength{\parskip}{6pt}\sloppy
\begin{document}
Abstract: Two invalid steps are identified in the core subcritical sticking argument, together with an admissible counterexample to the growing-SNR extension’s fluctuation mean. The fixed-SNR main limiting theorem is not disproved. Local concerns affect an iid-prior example and numerical parameter descriptions. All scientific judgments are pending. Original-page inspection and independent calculations support the scoped findings.\par
Author: AI-assisted technical review\par
Date:\par
Title: Initial scientific review: Fluctuations of the Largest Eigenvalues of Transformed Spiked Wigner Matrices\par
\# Paper identity and summary\par
This initial review concerns Aro Lee and Ji Oon Lee, Fluctuations of the Largest Eigenvalues of Transformed Spiked Wigner Matrices, ICML 2025, PMLR 267, paper 0193 (I25-43562), OpenReview xraEFvFHR4 [Retained latex: \textbackslash{}cite\{paper\}]. The unchanged staged proceedings PDF has 30 pages. References below use one-based PDF pages, which agree with printed page numbers. Exact source hashes are recorded in the reproducibility appendix and adjacent verification records.\par
Three Major findings are supported: an incorrect sign in a core sticking proposition, an incorrect edge count in its proof, and a wrong fluctuation mean in the growing-SNR extension. Local concerns affect an iid-prior example and numerical parameter descriptions. The fixed-SNR main limiting theorem is not disproved by these checks. All judgments remain pending human adjudication.\par
\# Scope and necessary background\par
The main assumptions, theorem statements, proof strategy, numerical examples, Taylor reduction, subcritical argument, and growing-SNR extension were inspected. Selected supercritical proof steps were cross-checked; this is not an independent reconstruction of every local-law estimate. No released implementation or raw simulation draws were available in the staged package. Separate supplementary-file availability remains unresolved.\par
For normalized noise \(Z_{ij}=\sqrt N W_{ij}\), the fixed-\(\lambda\) effective SNR is \[\lambda_e=\lambda\bigl(\mathbb E f'(Z_{12})\bigr)^2.\] The main theorem gives an outlier Gaussian limit for \(\lambda_e>1\) and a Tracy–Widom limit for \(\lambda_e<1\). The critical case is explicitly outside its scope. The reduced transformed matrix retains a secondary rank-one term involving \(f''\); its sign and scale matter below.\par
\# Problem 1: an invalid signed sticking assertion\par
\#\# Claim, assumptions, and dependency\par
Proposition 4.1 (p. 7) states, in paraphrase, \[0\leq \mu_1(H)-\mu_1(V)\leq N^{-3/4}\] with overwhelming probability in the subcritical regime. The proof on p. 24 declares the lower bound immediate from min–max. This proposition is a core bridge between the reduced transformed matrix and the noise matrix in the subcritical fluctuation proof.\par
Assumption 2.4 requires \(\mathbb E f'\geq0\), but not \(\mathbb E f''\geq0\). The temporary nonnegative-curvature assumption in Appendix B.3 simplifies that calculation and mentions the negative case separately; it is not a global restriction.\par
\#\# Decisive admissible counterexample\par
Take even \(N\), independent standard Gaussian \(Z_{ij}\) for \(i\leq j\), including the diagonal, and a deterministic balanced vector \(x_i\in\{-N^{-1/2},N^{-1/2}\}\). Then \[\|x\|=1,\qquad \sum_i x_i=\sum_i x_i^3=0.\] Choose fixed \(\lambda>0\) and \[f(z)=-\frac{z^2-1}{\sqrt2}.\] This polynomial has the required regularity and moments: \[\mathbb E f=0,\quad \mathbb E f^2=1,\quad
 \mathbb E f'=0,\quad \mathbb E f''=-\sqrt2,\quad \lambda_e=0<1.\] Writing \(u=N^{-1/2}(1,\ldots,1)^\top\), the retained term is \[H-V=-\frac{\lambda\sqrt N}{\sqrt2}x^2(x^2)^\top
     =-\frac{\lambda}{\sqrt{2N}}uu^\top.\] There is no Taylor remainder for this quadratic function. Thus min–max gives \(\mu_1(H)\leq\mu_1(V)\), in the opposite direction. Strict inequality holds almost surely: the independent entries have continuous laws, and equality would require the top eigenspace to be orthogonal to \(u\), an exceptional algebraic condition.\par
An independently written NumPy illustration with \(N=64\), \(\lambda=1\), and seed 193 returned \[\mu_1(V)=1.874218386760028,\qquad
 \mu_1(H)=1.8727533380544827.\] The analytic family, rather than this one draw, establishes the failure. The script and actual output are evidence E3–E4; the original proof page is E5.\par
\#\# Counterevidence, impact, and correction\par
A two-sided sticking estimate may remain true, and the fixed-\(\lambda\) Tracy–Widom conclusion is not refuted. The displayed proposition and its first proof step are nevertheless false under the stated hypotheses. Assessment: Major, because this invalid derivation lies directly on the main subcritical proof path.\par
A correction must prove the requisite two-sided estimate for both signs or explicitly restrict all affected statements. Reversing the inequality in this example is not a general repair. No complete replacement proof is certified.\par
The sign also affects Proposition 3.5’s ordered-spectrum wording. For \(f(z)=(-z^2+3z+1)/\sqrt{11}\), Gaussian noise, the balanced spike, and \(\lambda=4\), the nonzero eigenvalues of the deterministic perturbation are \(6/\sqrt{11}\) and \(-4/\sqrt{11N}\). The negative one is last in decreasing order, not second. The weaker assertion of at most two nonzero eigenvalues is compatible with the intended positive-outlier argument.\par
\# Problem 2: an incorrect edge count in the core sticking proof\par
Appendix C.1, p. 25, defines \[E_3=\{\mu_i(V):N^{-5/9}<\mu_1(V)-\mu_i(V)\leq N^{-2/9}\}\] and asserts \(|E_3|=O(N^{1/3})\). Equation (71) uses this cardinality to control the imaginary part of the unnormalized resolvent trace.\par
Set \(f(z)=z\), use normalized independent Gaussian entries and the balanced spike, and choose \(\lambda=1/4\). These are admissible and \(V=W\). At the right semicircle edge, \(\rho_{\rm sc}(2-s)\sim\sqrt{s}/\pi\). The eigenvalue rigidity invoked by the paper gives \[|E_3|\sim\frac{2N}{3\pi}
   \left\{(N^{-2/9})^{3/2}-(N^{-5/9})^{3/2}\right\}
   \asymp N^{2/3}.\] The \(O_\prec(N^{-2/3})\) edge displacement and rigidity errors do not change this leading count. Thus the printed \(N^{1/3}\) bound is false even in an ordinary admissible Wigner case.\par
With the corrected count, but otherwise retaining the displayed crude calculation, the contribution is controlled only at scale \[N^{2/3}N^{-5/6}N^{10/9}=N^{17/18}.\] This exceeds the \(N^{5/6}\) lower scale needed for the intended contradiction. The displayed argument therefore does not prove its conclusion. The original page and independent exponent calculation are E6 and E4.\par
A finer summation over edge distances can improve this crude bound substantially. The review does not claim that the actual trace reaches \(N^{17/18}\), that the limiting theorem is false, or that repair is impossible. Assessment: Major, for an invalid derivation in the main sticking proof. A full repaired summation with its probability and spectral-parameter control remains unverified.\par
\# Problem 3: a wrong fluctuation mean in the growing-SNR extension\par
\#\# An exact admissible quadratic construction\par
Appendix D, pp. 27–30, separately considers \(\mathbb E f'=0\) and \(\lambda=\lambda_0\sqrt N\). This problem concerns that extension, not the fixed-\(\lambda\) main theorem.\par
Let \(Z_{ij}\) be independent and uniform on \[\{-\sqrt{3/2},-\sqrt{1/2},\sqrt{1/2},\sqrt{3/2}\}.\] Let \(s_i\) be independent Rademacher variables independent of the noise, and set \[x_i=s_i/\sqrt N,\qquad f(z)=2(z^2-1),\qquad \lambda_0=1.\] The noise has mean zero, variance one, fourth moment \(5/4\), and bounded moments of every order. The transformation satisfies \[\mathbb E f=0,\quad \mathbb E f^2=1,\quad
 \mathbb E f'=0,\quad \mathbb E f''=4.\] The spike has norm exactly one and satisfies all delocalization and sum conditions with overwhelming probability. Indeed \(\sum_i x_i^3=N^{-1}\sum_i x_i=O_\prec(N^{-1})\). Remark 2.8 gives the effective SNR directly as \[\widetilde\lambda_e=\tfrac14N^2(4)^2\|x^2\|^4=4.\] No reliance on the ambiguous fourth-moment notation on p. 30 is needed.\par
The exact transformed matrix is \[\widetilde M=V+2uu^\top,\qquad
 V_{ij}=\frac{2(Z_{ij}^2-1)}{\sqrt N}
       +4Z_{ij}s_is_jN^{-3/4}.\] Symmetry of \(Z_{ij}\) makes the conditional joint distribution of \(V\) independent of the signs; its upper-triangle entries remain iid. With \(d=N^{-1/2}\), \[\mathbb E V_{ij}=0,\qquad
 \mathbb E V_{ij}^2=\frac{1+16d}{N}.\] The true variance correction coefficient is 16.\par
\#\# Variance coefficient and normalized-spike errors\par
Equation (75), p. 28, literally begins its variance coefficient with \(\mathbb E[f_{12}^2]\), without a prime, as confirmed on the rendered original page E7. Its printed coefficient equals 1 in this example. Substituting the likely intended \(\mathbb E[(f'_{12})^2]\) gives 16.\par
Even after that charitable correction, Eq. (79) has a distinct centering error. Define \[a_N=\sqrt{1+16d},\quad W_N=V/a_N,\quad \theta_N=2/a_N.\] Then \(\widetilde M=a_N(W_N+\theta_Nuu^\top)\): the normalized spike changes with the noise scale. Exact algebra gives \[\theta_N+\theta_N^{-1}=\frac52-12d+O(d^2),\qquad
 a_N(\theta_N+\theta_N^{-1})=\frac52+8d.\] Multiplying only the fixed center \(5/2\), as in the paper, yields instead \(5/2+20d+O(d^2)\). The omitted spike change is precisely at the fluctuation scale.\par
\#\# The discrepancy survives in the limiting law\par
The normalized entries of \(W_N\) are uniformly bounded, centered, and variance one. Their third moment is \[\gamma_{3,N}=\frac{24d}{(1+16d)^{3/2}}\to0.\] The fixed-spike delocalized-outlier theorem of Renfrew and Soshnikov [Retained latex: \textbackslash{}cite\{rs\}], Theorem 1.3 under (1.1)–(1.6), applies to \(B_N=W_N+2uu^\top\). Its third-moment correction vanishes here, and its variance conversion gives \[\sqrt N\left(\lambda_{\max}(B_N)-\frac52\right)
       \Rightarrow \mathcal{N}(0,3/2).\] The external theorem’s fixed-spike requirement is respected; no unverified varying-spike extension is assumed.\par
For the remaining change, put \(q_N(z)=u^\top(z-W_N)^{-1}u\). Its outliers satisfy \(q_N(\lambda(\theta))=1/\theta\). The quadratic-form semicircle limit, obtained from Theorem 1.6 of [Retained latex: \textbackslash{}cite\{rs\}] using a smooth cutoff of the squared resolvent, gives \(q'_N(\zeta_N)\to-1/3\) for intermediate points \(\zeta_N\to5/2\); the bulk remains separated. Since \[\theta_N^{-1}-\tfrac12=4d+O(d^2),\] the mean-value theorem gives \[\lambda(\theta_N)-\lambda(2)=-12d+o_p(d).\] Multiplying back by \(a_N=1+8d+O(d^2)\) therefore proves \[\boxed{\sqrt N\left(\lambda_{\max}(\widetilde M)-\frac52\right)
          \Rightarrow \mathcal{N}(8,3/2).}\] The parent checked the primary theorem hypotheses and derivative argument. An independent SymPy calculation verified the exact moments and center expansions (E3–E4). The external-source check is preserved in E12. This is an analytical conclusion supported by a primary theorem, not a claim of empirical reproduction.\par
Equation (79) predicts mean \(5/4\) with the literal Eq. (75), or 20 after the missing-prime correction. The correct mean is 8. Assessment: Major: this extension’s quantitative conclusion is false in a fully admissible case. The verified correction applies to this example; a general replacement theorem has not been proved.\par
\# Problem 4: the iid-prior example omits a third-moment condition\par
The remark following Assumption 2.3 (p. 4) supplies iid mean-zero, unit-variance, finite-moment coordinates as an example. The assumption additionally requires \(\sum_i x_i^3=O(N^{-1+\varepsilon})\) for every \(\varepsilon>0\).\par
Let \(Y=\sqrt2\) with probability \(1/3\), and \(Y=-1/\sqrt2\) with probability \(2/3\). Then \[\mathbb EY=0,\qquad\mathbb EY^2=1,\qquad\mathbb EY^3=1/\sqrt2.\] For \(x_i=Y_i/\sqrt N\), the law of large numbers gives \[\sqrt N\sum_i x_i^3\to1/\sqrt2.\] The same holds after exact normalization \(x_i=Y_i/\sqrt{\sum_jY_j^2}\). Hence the cubic sum has order \(N^{-1/2}\), violating the claimed condition when \(\varepsilon<1/2\).\par
Assessment: Minor, a local scope error. The theorem itself explicitly assumes the stronger condition; the symmetric Rademacher simulation prior satisfies it. A zero-third-moment restriction, with appropriate normalization and moment control, repairs this example. The review does not extend the theorem to inadmissible asymmetric priors.\par
\# Problem 5: inconsistent numerical parameter descriptions\par
For the Gaussian example on pp. 9 and 13, \[f(z)=\frac{z^2+3z-1}{\sqrt{11}},\qquad\lambda_e=\frac9{11}\lambda.\] Thus \(\lambda=2.5\) gives \(\lambda_e=45/22\approx2.04545\), not the printed 1.294. The prose’s \(\lambda=0.1\) gives approximately 0.081818, not 0.350. Figure 3 instead captions the subcritical experiment with \(\lambda=0.15\), giving approximately 0.122727. These discrepancies change the stated experiment and, in the supercritical case, the theoretical Gaussian variance.\par
There is substantial counterevidence against inferring that the plots used the wrong theory: the visually inspected upper Gaussian curve peaks near 0.4, consistent with the correct variance \(46/45\). Without code or raw draws, the actual generating settings remain unresolved. Assessment: Minor, for a reproducibility and calculation-description defect, not a demonstrated failed empirical fit.\par
For completeness, Eq. (25), p. 13, has another local normalization error. The stated mixture \((G+2R)/\sqrt5\) needs prefactor \(\sqrt5/(2\sqrt{2\pi})\), whereas the printed \(5/(2\sqrt{2\pi})\) integrates to \(\sqrt5\). A constant cancels in the score \(-p'/p\), and sampling from the stated mixture can remain correct. This typographical observation is not a separate substantive finding. Source pages 9 and 13 and the independent arithmetic are E9–E10 and E4.\par
\# Overall impact, positive checks, and remaining scope\par
The invalid sign and edge count affect the displayed derivation of core subcritical sticking. They do not establish a false Tracy–Widom limit. The Appendix D counterexample does establish a wrong limiting mean for the stated growing-SNR scaling. Its appendix location does not diminish that quantitative defect, but its different scaling limits its effect on the main theorem.\par
The specialization \(f(z)=z\), \(\lambda=4\) agrees with ordinary spiked Wigner behavior: center \(5/2\), variance \(3/2\). The main Taylor reduction correctly retains the quadratic rank term; its remaining \(O_\prec(N^{-1})\) error is smaller than both relevant fluctuation scales. Polynomial \(C^3\) control supports the remainder rate.\par
All supercritical comparison estimates, every histogram, and a repaired subcritical proof were not independently reproduced. No author correction or adoption was established in the inspected sources. No inference of first discovery is made.\par
\# Conclusion and state table\par
Item | Assessment | Repair verification | Author adoption | Remaining work\par
P1 | Major | Counterexample verified; general repair absent | Not requested | Two-sided sticking proof\par
P2 | Major | False count verified; sharper proof absent | Not requested | Valid trace summation\par
P3 | Major | Exact example corrected; general repair absent | Not requested | General extension\par
P4 | Minor | Local hypothesis correction checked | Not requested | Correct prior scope\par
P5 | Minor | Arithmetic checked | Not requested | Actual simulation settings\par
The report-level assessment is Major. Scientific adjudication is pending and human reviews are empty. The original scientific review scope is preserved; focused acceptance does not constitute human adjudication.\par
\# Reproducible evidence and execution record\par
The source PDF SHA-256 is\par
964ce2fe0737106ccd7a215e3ceb67876a0e5ea7cc432874235ef561101b19d5.\par
The manifest, extracted text, independent scripts/results, inspected page images, and report source are hash-bound in the adjacent verification records. The parent actually inspected rendered source pages 9, 13, 24, 25, 28, and 30. Existing NumPy and SymPy were used for independent checks; no author code, model, or dataset was executed or downloaded.\par
Three native auxiliary agents assisted with theorem assumptions, numerical consistency, and Appendix D. Their exposed inherited identity was GPT-6 Astra Pro; no backend version ID was exposed. The parent integrated their work, independently read decisive original passages and primary external results, and owns the judgments. No separate generating conversation was created.\par
The frozen mathematical-journal template preamble was retained. The initial Windows review had no available TeX compiler. During focused recovery, the final TeX was compiled with the existing isolated runtime’s pdfLaTeX and all resulting report pages were inspected. Compiler, dependency, PDF, and render hashes are recorded in provenance. No software installation or unrelated cache import was used.\par
9 A. Lee and J. O. Lee, Fluctuations of the Largest Eigenvalues of Transformed Spiked Wigner Matrices, ICML, PMLR 267 (2025). https://proceedings.mlr.press/v267/lee25t.html. Reviewed version: staged proceedings PDF specified above. D. Renfrew and A. Soshnikov, On Finite Rank Deformations of Wigner Matrices II: Delocalized Perturbations, arXiv:1203.5130v3, 22 October 2012. https://arxiv.org/pdf/1203.5130v3. Consulted online on 24 October 2026: assumptions (1.1)–(1.6), fixed-spike restriction, Theorems 1.3 and 1.6, and derivative identity (1.24). The consulted version is now preserved at ; its actual hash is recorded in provenance.\par
\end{document}
